\documentclass[11pt,a4paper]{article}

% ------------------------------------------------------------------ %
% Packages
% ------------------------------------------------------------------ %
\usepackage[margin=2.2cm]{geometry}
\usepackage{amsmath,amssymb,mathtools}
\usepackage{booktabs}
\usepackage{longtable}
\usepackage{array}
\usepackage{xcolor}
\usepackage{listings}
\usepackage{siunitx}
\usepackage{microtype}
\usepackage{enumitem}
\usepackage{algorithm}
\usepackage{algpseudocode}
\usepackage[hidelinks]{hyperref}

% ------------------------------------------------------------------ %
% Listing style for parameter snippets
% ------------------------------------------------------------------ %
\lstdefinestyle{cfg}{
  basicstyle=\ttfamily\footnotesize,
  breaklines=true,
  columns=fullflexible,
  frame=single,
  rulecolor=\color{black!30},
  showstringspaces=false,
}
\lstset{style=cfg}

% Tighter algorithmic display
\algrenewcommand\algorithmicrequire{\textbf{Input:}}
\algrenewcommand\algorithmicensure{\textbf{Output:}}
\newcommand{\pyconst}[1]{\textcolor{black!60!blue}{\texttt{\footnotesize #1}}}
\newcommand{\dictkey}[1]{\textcolor{black!60!green}{\texttt{\footnotesize #1}}}

% ------------------------------------------------------------------ %
% Convenience macros
% ------------------------------------------------------------------ %
\newcommand{\Vm}{V_{m}}
\newcommand{\Vmsub}[1]{V_{m,#1}}
\newcommand{\Iion}{I_{\text{ion}}}
\newcommand{\Iionsub}[1]{I_{\text{ion},#1}}
\newcommand{\Iext}{I_{\text{ext}}}
\newcommand{\Dten}{\mathbf{D}}
\newcommand{\grad}{\nabla}
\newcommand{\divg}{\nabla\!\cdot}
\newcommand{\Aimp}{\mathsf{A}}
\newcommand{\Bimp}{\mathsf{B}}
\newcommand{\Mmass}{\mathsf{M}}
\newcommand{\Kstif}{\mathsf{K}}
\newcommand{\Lop}{\mathsf{L}}
\newcommand{\code}[1]{\texttt{\small #1}}

\title{Solver Specification:\\
       \texorpdfstring{\texttt{highOrderElectroActivationFoamImplicitJfNK}}{highOrderElectroActivationFoamImplicitJfNK}\\
       \large A high-order, implicit, Jacobian-free Newton--Krylov
       monodomain solver in OpenFOAM with an Eigen linear-algebra back-end}
\author{Pablo}
\date{\today}

\begin{document}
\maketitle

\begin{abstract}
This report specifies the algorithms implemented in the OpenFOAM solver
\code{highOrderElectroActivationFoamImplicitJfNK}. The solver
advances the monodomain equations of cardiac electrophysiology coupled
to the ten Tusscher--Noble--Noble--Panfilov (TNNP, 2004) ionic cell
model. The diffusion operator is discretised with the cell-centred
finite-volume framework of OpenFOAM, optionally upgraded to a Local
Radial-Extrapolation (LRE) high-order reconstruction; the ionic source
is evaluated either at cell centres or at LRE cell quadrature points;
time integration follows a $\theta$-scheme (backward Euler or
Crank--Nicolson); the resulting nonlinear coupled system is solved by
a Jacobian-Free Newton--Krylov (JFNK) method with Armijo backtracking
line search, linear extrapolation initial guess and a physical-bracket
clamp on the input to the ionic ODE; the inner linear system is solved
by an in-house restarted GMRES with modified Gram--Schmidt
orthogonalisation built on top of the Eigen sparse linear algebra
library; the per-cell ionic ODE uses OpenFOAM's adaptive RKF45
integrator with configurable tolerances; activation times are tracked
by linear interpolation across the first threshold crossing and
benchmark points are sampled by inverse-distance weighting. All
algorithmic parameters are exposed through the
\code{constant/spatialIntegrationProperties} and
\code{constant/timeIntegrationProperties} dictionaries, and the
companion Python driver \code{run\_convergence.py} generates these
dictionaries from a single configuration file. This solver is
algorithmically identical to its PETSc-based sibling
\code{highOrderElectroActivationFoamImplicitJfNK\_PETSc}; the only
difference is the linear-algebra library (Eigen vs.\ PETSc).
\end{abstract}

\tableofcontents
\bigskip

% ================================================================== %
\section{Notation}
% ================================================================== %

\begin{center}
\begin{tabular}{@{}lll@{}}
\toprule
Symbol & Meaning & Units \\
\midrule
$\Omega \subset \mathbb{R}^{d}$ & Spatial domain ($d \in \{2,3\}$) & \si{\meter} \\
$\partial\Omega$ & Insulated boundary (zero-flux Neumann) & --- \\
$t \in (0, T]$ & Time & \si{\second} \\
$\Vm(\mathbf{x}, t)$ & Transmembrane voltage & \si{\volt} \\
$\mathbf{s}(\mathbf{x}, t) \in \mathbb{R}^{17}$ & TNNP ionic states & mixed \\
$\Iion(\Vm, \mathbf{s})$ & Total ionic current density & \si{\volt\per\second} \\
$\Iext(\mathbf{x}, t)$ & External stimulus current density & \si{\ampere\per\meter\cubed} \\
$\chi$ & Membrane surface-to-volume ratio & \si{\per\meter} \\
$C_{m}$ & Membrane capacitance per unit area & \si{\farad\per\meter\squared} \\
$\Dten(\mathbf{x})$ & Bulk conductivity tensor & \si{\siemens\per\meter} \\
$\theta \in \{0.5, 1\}$ & Time-stepping parameter (CN or BE) & --- \\
$\Delta t$ & Outer time-step length & \si{\second} \\
$N_{c}$ & Number of mesh cells & --- \\
$\mathsf{V}^{n} \in \mathbb{R}^{N_{c}}$ & Vector of cell-averaged $\Vm$ at time $t^{n}$ & \si{\volt} \\
$\Mmass, \Kstif, \Lop$ & Mass, stiffness, scaled diffusion matrix & --- \\
$\Aimp, \Bimp$ & Implicit LHS / RHS operators of the $\theta$-scheme & --- \\
$\mathsf{F}(\mathsf{x})$ & Discrete nonlinear residual to be zeroed & --- \\
\bottomrule
\end{tabular}
\end{center}

% ================================================================== %
\section{Governing equations}
% ================================================================== %

\subsection{Monodomain PDE}

The bulk voltage is propagated by the standard monodomain
reaction--diffusion equation~\cite{ttp2004,niederer2012}:
\begin{equation}
  \chi \, C_{m} \, \frac{\partial \Vm}{\partial t}
   \;=\;
   \divg\!\bigl( \Dten \grad \Vm \bigr)
   - \Iion(\Vm, \mathbf{s})
   + \Iext
   \quad \text{in } \Omega \times (0, T],
   \label{eq:monodomain}
\end{equation}
subject to insulated boundaries
$\mathbf{n} \cdot ( \Dten \grad \Vm ) = 0$ on
$\partial \Omega \times (0, T]$, and initial condition
$\Vm(\mathbf{x}, 0) = -\SI{84}{\milli\volt}$. The conductivity tensor
$\Dten$ is symmetric positive-definite and either uniform or read as a
\code{volTensorField} from disk.

Dividing~\eqref{eq:monodomain} by $\chi C_{m}$ gives the form actually
used inside the solver,
\begin{equation}
  \frac{\partial \Vm}{\partial t}
   \;=\;
   \mathcal{L} \Vm
   \;+\;
   s(\Vm, \mathbf{s}),
   \qquad
   \mathcal{L} \;=\; \frac{1}{\chi C_{m}}\,\divg(\Dten \grad \cdot),
   \quad
   s \;=\; -\frac{\Iion}{\chi C_{m}} + \frac{\Iext}{\chi C_{m}}.
   \label{eq:semi-discrete}
\end{equation}

\subsection{Ionic ODE}

At every spatial sampling point the TNNP state vector
$\mathbf{s} \in \mathbb{R}^{17}$ evolves according to
\begin{equation}
  \frac{\mathrm{d}\mathbf{s}}{\mathrm{d}t}
   \;=\;
   \mathbf{f}_{\text{TNNP}}\!\bigl(\Vm, \mathbf{s}; \texttt{tissue}\bigr),
   \qquad \mathbf{s}(0) = \mathbf{s}_{0},
   \label{eq:tnnp}
\end{equation}
where \code{tissue} selects \code{epicardialCells}, \code{mCells} or
\code{endocardialCells}. The 17 tracked states are
\[
  \mathbf{s} = (V,\ K_{i},\ Na_{i},\ Ca_{i},\ Xr_{1},\ Xr_{2},\ Xs,\ m,\ h,\ j,\ d,\ f,\ fCa,\ s,\ r,\ g,\ Ca_{SR}).
\]
The total ionic current returned to the PDE is the algebraic
\code{Iion\_cm} of TNNP~\cite{ttp2004}.

\subsection{External stimulus}

The external current is a time- and region-windowed step:
\begin{equation}
  \Iext(\mathbf{x}, t) \;=\;
  \begin{cases}
    I_{0} & \mathbf{x} \in \Omega_{\text{stim}}^{(b)},
            \;\; t \in [t_{0}^{(b)},\, t_{0}^{(b)} + \tau],
            \;\; b = 1, \dots, B, \\
    0     & \text{otherwise},
  \end{cases}
\end{equation}
with $I_{0} = \SI{50000}{\ampere\per\meter\cubed}$ and
$\tau = \SI{2}{\milli\second}$ as Niederer-benchmark defaults.

% ================================================================== %
\section{Time discretisation: the $\theta$-scheme}
% ================================================================== %

Time is advanced from $t^{n}$ to $t^{n+1} = t^{n} + \Delta t$ by
\begin{equation}
  \frac{\mathsf{V}^{n+1} - \mathsf{V}^{n}}{\Delta t}
   \;=\;
   \theta\, \Lop\, \mathsf{V}^{n+1}
   + (1-\theta)\, \Lop\, \mathsf{V}^{n}
   + \theta\, \mathbf{s}^{n+1}
   + (1-\theta)\, \mathbf{s}^{n}.
   \label{eq:theta}
\end{equation}

\begin{center}
\begin{tabular}{@{}lcl@{}}
\toprule
\code{implicitScheme} & $\theta$ & Scheme \\
\midrule
\code{backwardEuler} & $1$   & First-order, L-stable \\
\code{crankNicolson} & $1/2$ & Second-order, A-stable \\
\bottomrule
\end{tabular}
\end{center}

After multiplication by $\Delta t^{-1}$ the implementation forms
\begin{equation}
  \bigl[\Mmass/\Delta t \,-\, \theta\, \Lop\bigr]\, \mathsf{V}^{n+1}
   \;=\;
   \bigl[\Mmass/\Delta t \,+\, (1-\theta)\, \Lop\bigr]\, \mathsf{V}^{n}
   \;+\; \theta\, \mathbf{s}^{n+1}
   \;+\; (1-\theta)\, \mathbf{s}^{n},
   \label{eq:implicit-system}
\end{equation}
with the two Eigen \code{SparseMatrix<scalar,RowMajor>} operators
\begin{equation}
  \Aimp \;=\; \Mmass/\Delta t \,-\, \theta\, \Lop,
   \qquad
  \Bimp \;=\; \Mmass/\Delta t \,+\, (1-\theta)\, \Lop.
   \label{eq:AB}
\end{equation}
$\Aimp$ and $\Bimp$ are rebuilt at every time step by copying $\Mmass$,
scaling by $\Delta t^{-1}$ and adding $\pm \theta \Lop$. When
$\theta = 1$ (BE) the explicit half of the diffusion vanishes and
$\Bimp = \Mmass/\Delta t$. The scaled diffusion is
$\Lop = \Kstif / (\chi C_{m})$.

% ================================================================== %
\section{Spatial discretisation}
% ================================================================== %

\subsection{Standard orthogonal FV Laplacian}

\code{assembleStandardOrthogonalStiffnessMatrix} computes the classical
5/7-point Laplacian. Per internal face $f$ with face area $\mathbf{S}_{f}$,
\begin{equation}
  a_{f} \;=\;
   \frac{ \lvert \mathbf{S}_{f} \rvert\,
          \bigl( \mathbf{n}_{f} \cdot \Dten_{f} \cdot \hat{\mathbf{d}}_{f} \bigr) }
        { \lvert \mathbf{d}_{f} \rvert },
   \quad
   \Dten_{f} = \tfrac{1}{2}\,(\Dten_{\text{own}} + \Dten_{\text{nei}}),
\end{equation}
contributing $-a_{f}/V_{\text{own}}$ to $\Kstif[\text{own},\text{own}]$,
$+a_{f}/V_{\text{own}}$ to $\Kstif[\text{own},\text{nei}]$, and the
transposed pair to the neighbour row. \code{zeroGradient}/\code{empty}
patches contribute zero flux. \code{fixedValue} patches contribute to
the owner diagonal only.

\subsection{High-order LRE diffusion}

When \code{useHighOrder\_Vm = true} and $d>1$,
\code{assembleHighOrderStiffnessMatrix} integrates the flux on a face
quadrature with reconstructed gradients:
\begin{equation}
  \Phi_{f} \;=\;
   \sum_{q \in \mathcal{Q}_{f}}
       w_{q}\, \lvert \mathbf{S}_{f} \rvert\,
       \mathbf{n}_{f} \cdot \Dten_{f} \cdot
       \Bigl(
         \sum_{j \in \mathcal{S}_{f}}
            \mathbf{g}_{f,q,j}\, \mathsf{V}_{j}
       \Bigr),
   \label{eq:HO-flux}
\end{equation}
where $\mathcal{S}_{f}$ is the LRE face stencil ($\sim 15$--$50$ cells)
and $\mathbf{g}_{f,q,j} = $~\code{QRGradFaceGPCoeffs}. The
contribution of each $\mathsf{V}_{j}$ is added to row \code{own} with
weight $+\Phi_{f,j}/V_{\text{own}}$ and to row \code{nei} with weight
$-\Phi_{f,j}/V_{\text{nei}}$, yielding a negative semi-definite
Laplacian.

\subsection{Mass matrices}

Two options:
\paragraph{Lumped (\code{massMatrix lumped}).}
\code{assembleDiagonalMassMatrix} writes $\Mmass = I$.

\paragraph{Consistent HO (\code{massMatrix consistent}).}
\code{assembleConsistentMassMatrixHO} fills
\begin{equation}
  \Mmass[i, j] \;=\;
   \sum_{q \in \mathcal{Q}_{i}}
       \frac{w_{q}}{\sum_{q'} w_{q'}}\,
       \phi_{j}(\mathbf{x}_{q}),
   \quad
   \phi_{j}(\mathbf{x}_{q})
    = \delta_{ij}
    + \mathbf{g}_{i,j} \cdot \mathbf{d}_{q}
    + \tfrac{1}{2}\, \mathbf{H}_{i,j} \!:\! (\mathbf{d}_{q} \otimes \mathbf{d}_{q})
    + \tfrac{1}{6}\, \mathbf{T}^{(3)}_{i,j} \!\vdots\!
        (\mathbf{d}_{q} \otimes \mathbf{d}_{q} \otimes \mathbf{d}_{q}),
   \label{eq:HO-mass}
\end{equation}
with $\mathbf{d}_{q} = \mathbf{x}_{q} - \mathbf{C}_{i}$. A safety guard
silently downgrades to the lumped form when
\code{useHighOrder\_Vm = false} or $d \le 1$.

\subsection{High-order reconstruction at ionic quadrature points}

If \code{useHighOrder\_Iion = true}, the ionic ODE is integrated at
cell quadrature points. Vm is reconstructed there by the same Taylor
expansion as in~\eqref{eq:HO-mass}:
\begin{equation}
  \Vm(\mathbf{x}_{q})
   \;\approx\;
   \Vmsub{c} + (\grad \Vm)_{c} \cdot \mathbf{d}_{q}
   + \tfrac{1}{2}\, \mathbf{H}_{c} \!:\! (\mathbf{d}_{q} \otimes \mathbf{d}_{q})
   + \tfrac{1}{6}\, \mathbf{T}^{(3)}_{c} \!\vdots\!
       (\mathbf{d}_{q} \otimes \mathbf{d}_{q} \otimes \mathbf{d}_{q}),
   \label{eq:HO-vm}
\end{equation}
implemented in \code{reconstructVmAtIionIntegrationPoints}.

\subsection{High-order cell averaging of $\Iion$}

\code{averageIionIntegrationPointsToCells} contracts the per-QP $\Iion$
back to cells:
\begin{equation}
  \Iionsub{i} \;=\;
   \frac{\sum_{q \in \mathcal{Q}_{i}} w_{q}\, \Iionsub{q}}
        {\sum_{q \in \mathcal{Q}_{i}} w_{q}}.
\end{equation}

% ================================================================== %
\section{Nonlinear solver}
% ================================================================== %

The fully discrete equation is nonlinear in $\mathsf{V}^{n+1}$. We
define
\begin{equation}
  \mathsf{F}(\mathsf{x})
   \;=\;
   \Aimp\, \mathsf{x}
   \;-\;
   \mathsf{r}\!\bigl( \Iion(\mathsf{x}) \bigr),
   \qquad
   \mathsf{x} = \mathsf{V}^{n+1},
   \label{eq:Fdef}
\end{equation}
with
$\mathsf{r}(\Iion) = \Bimp\, \mathsf{V}^{n} +
\theta\, \mathbf{s}^{n+1} + (1-\theta)\, \mathbf{s}^{n}$.

\subsection{Jacobian-Free Newton--Krylov}

Newton iterates $\mathsf{J}(\mathsf{x}_{k})\,\boldsymbol{\delta}_{k} =
-\mathsf{F}(\mathsf{x}_{k})$, with $\mathsf{J}$ never assembled.
Each matrix--vector product needed by GMRES is approximated by a
forward finite difference~\cite{knoll2004}:
\begin{equation}
  \mathsf{J}(\mathsf{x}_{k})\, \mathbf{v}
   \;\approx\;
   \frac{\mathsf{F}(\mathsf{x}_{k} + \varepsilon\, \mathbf{v})
         - \mathsf{F}(\mathsf{x}_{k})}{\varepsilon},
   \qquad
   \varepsilon \;=\; \varepsilon_{0}\,
                     \frac{1 + \lVert \mathsf{x}_{k} \rVert_{2}}
                          {\lVert \mathbf{v} \rVert_{2}},
   \label{eq:fdjac}
\end{equation}
with $\varepsilon_{0} = $~\code{jfnkEpsilon}. The well-known
optimisation is $\varepsilon_{0} \approx \sqrt{\eta}$ where $\eta$ is
the noise floor of $\mathsf{F}$ (here driven by the RKF45 \code{relTol}).

\subsection{In-house restarted GMRES}

\code{solveGMRES} is an in-house implementation of GMRES(m) with
modified Gram--Schmidt orthogonalisation and outer restarts. Each
inner iteration:
\begin{enumerate}[topsep=0pt,itemsep=2pt]
  \item builds one Krylov direction via \code{matVec};
  \item orthogonalises against the running basis $V$ with MGS, updating
        the upper-Hessenberg matrix $H$;
  \item solves the small least-squares problem
        $\min \lVert g - H_{j} y \rVert$ with Eigen's
        column-pivoted Householder QR;
  \item estimates the relative residual exactly as
        $\lVert g - H_{j} y \rVert / \lVert b \rVert$ (no extra
        matVec).
\end{enumerate}
After \code{jfnkMaxKrylovIterations} inner steps the Krylov subspace
is discarded and the true residual $r = b - A x$ is recomputed
(one extra matVec) before the next restart, repeated up to
\code{jfnkMaxRestarts} times.

\subsection{Armijo backtracking line search}

The unconditional update $\mathsf{x}_{k+1} = \mathsf{x}_{k} +
\alpha_{0}\, \boldsymbol{\delta}_{k}$ is wrapped in Armijo
backtracking~\cite{noceWright}:
\begin{equation}
  \alpha_{k} \;=\; \alpha_{0} \cdot 2^{-\ell_{k}},
   \quad
   \ell_{k} \;=\; \min\Bigl\{\, \ell \ge 0 \;\big|\;
       \lVert \mathsf{F}(\mathsf{x}_{k} + 2^{-\ell}\alpha_{0}
                          \boldsymbol{\delta}_{k}) \rVert_{2}^{2}
       \le \bigl(1 - 2c\, 2^{-\ell}\alpha_{0}\bigr)\,
           \lVert \mathsf{F}(\mathsf{x}_{k}) \rVert_{2}^{2}
       \Bigr\},
   \label{eq:armijo}
\end{equation}
with $c = $~\code{jfnkArmijoC}. Backtracking stops on the first success
or after \code{jfnkLineSearchMaxIter} trials; an absolute floor
\code{jfnkLineSearchAlphaMin} below which $\alpha$ is forced and the
current trial accepted protects against infinite loops. Disable with
\code{jfnkLineSearch~false}.

\subsection{Linear extrapolation initial guess}

\code{jfnkInitGuessOrder} selects between
\begin{align}
  \mathsf{x}_{0} &\;=\; \mathsf{V}^{n}
       & \text{(\code{jfnkInitGuessOrder = 0})}, \\
  \mathsf{x}_{0} &\;=\; \mathsf{V}^{n} + \frac{\Delta t}{\Delta t_{\text{prev}}}
        \bigl( \mathsf{V}^{n} - \mathsf{V}^{n-1} \bigr)
       & \text{(\code{jfnkInitGuessOrder = 1}, default).}
   \label{eq:extrap}
\end{align}
Each cell is clamped to $[\,$\code{jfnkInitGuessVmMin},
\code{jfnkInitGuessVmMax}$\,]$. On the first time step no
$\mathsf{V}^{n-1}$ exists; the solver falls back to order 0.

\subsection{Protective Vm clamp on ODE input}

When \code{jfnkClampODEInput = true} (default), the Vm vector passed to
\code{TNNP::solveODE} inside every call to \code{evaluateNonlinearFields}
is clamped to $[\,$\code{jfnkInitGuessVmMin}, \code{jfnkInitGuessVmMax}$\,]$
\emph{in the working buffer only}: \code{VmIntegrationPoints} (HO-Iion
path) or a local \code{scalarField} copy (otherwise). The persistent
fields are never mutated.

\subsection{Convergence criteria}

Newton converges when all of the following hold:
\begin{equation*}
\begin{aligned}
  \lVert \mathsf{F}(\mathsf{x}_{k}) \rVert_{2} / \lVert \mathsf{x}_{k} \rVert_{2}
     &\;\le\; \texttt{nonlinearTolerance}, \\
  \lVert \mathsf{V}_{k} - \mathsf{V}_{k-1} \rVert_{2} /
     \lVert \mathsf{V}_{k} \rVert_{2}
     &\;\le\; \texttt{nonlinearVmTolerance}, \\
  \lVert \Iionsub{k} - \Iionsub{k-1} \rVert_{2} /
     \lVert \Iionsub{k} \rVert_{2}
     &\;\le\; \texttt{nonlinearIionTolerance}, \\
  \text{(opt.)}\quad
  \max_{i}\;
  \lVert \mathbf{s}_{k}^{(i)} - \mathbf{s}_{k-1}^{(i)} \rVert_{2} /
     \lVert \mathbf{s}_{k}^{(i)} \rVert_{2}
     &\;\le\; \texttt{nonlinearStatesTolerance},
\end{aligned}
\end{equation*}
plus $k+1 \ge $~\code{minNonlinearIterations}. The last criterion is
gated by \code{nonlinearRequireStatesConvergence}. The coupled residual
$\lVert \mathsf{F} \rVert / \lVert \mathsf{x} \rVert$ has a noise floor
$\sim \eta = $~\code{relTol} (RKF45); tolerances below this floor
cannot be achieved.

\subsection{Rollback on non-convergence}

If Newton exits with the iteration budget reached and the convergence
flag still false, \code{nonlinearAcceptUnconverged} decides:
\begin{itemize}[topsep=0pt]
  \item \code{false} (default): restore $\Vm, \Iion, \mathbf{s}$ to
        $t^{n}$ values, warn, and continue;
  \item \code{true}: commit the unconverged iterate.
\end{itemize}

\subsection{Alternative nonlinear methods}

Selectable by \code{nonlinearMethod}:
\begin{description}[style=nextline]
  \item[\code{picard}] Fixed-point with lagged $\Iion$ source.
  \item[\code{diagonalIion}] Picard with the diagonal of
        $\partial \Iion / \partial \Vm$ (one-sided FD step
        \code{diagonalIionEpsilon}) folded into $\Aimp$.
\end{description}

% ================================================================== %
\section{Linear solver: Eigen back-end}
% ================================================================== %

The wrapper \code{solveSparseSystem} drives either of two Eigen sparse
solvers on the explicit operator $\Aimp$ (Picard / diagonal paths):

\begin{center}
\begin{tabular}{@{}lll@{}}
\toprule
\code{linearSolver} key & Eigen type & Type \\
\midrule
\code{SparseLU}  & \code{Eigen::SparseLU} & Direct LU \\
\code{BiCGSTAB}  & \code{Eigen::BiCGSTAB} + \code{IncompleteLUT} & Iterative \\
\bottomrule
\end{tabular}
\end{center}

Both share the \code{tolerance} and \code{maxIterations} dictionary
keys (the former is ignored by \code{SparseLU} since it is direct).
The matrix-free JFNK branch never instantiates these wrappers; instead
it calls the in-house GMRES (Section~5.2) on the FD Jacobian shell.

% ================================================================== %
\section{Ionic ODE integration}
% ================================================================== %

For each cell or cell QP, the TNNP ODE~\eqref{eq:tnnp} is advanced by
the OpenFOAM adaptive integrator selected via the \code{solver} key.
Default \code{RKF45}~\cite{hairer2}. Keys read from
\code{constant/timeIntegrationProperties}:

\begin{center}
\begin{tabular}{@{}lll@{}}
\toprule
Key & Meaning & Default \\
\midrule
\code{solver} & Concrete ODE integrator & \code{RKF45} \\
\code{absTol} & Absolute tolerance per state & $10^{-9}$ \\
\code{relTol} & Relative tolerance per state & $10^{-6}$ \\
\code{maxSteps} & Cap on sub-steps per outer call & $10^{4}$ \\
\code{initialODEStep} & Initial sub-step (ms) & $10^{-5}$ \\
\bottomrule
\end{tabular}
\end{center}

\code{relTol} drives the noise floor of the matrix-free Jacobian:
$\eta_{\mathsf{F}} \sim $~\code{relTol}. Tightening \code{relTol} buys
a lower attainable outer tolerance at the cost of more RKF45 sub-steps
per Newton iteration.

% ================================================================== %
\section{TNNP numerical protection}
% ================================================================== %

The master switch \code{TNNPNumericalProtection} (default \code{false})
gates a guard \code{protectState(state, t, qp, phase)} invoked at every
external entry into the TNNP kernel (\code{calculateCurrent},
\code{solveODE\_start}, \code{solveODE\_end}, \code{derivatives}).

\begin{center}
\begin{tabular}{@{}lll@{}}
\toprule
Key & Type & Default \\
\midrule
\code{TNNPNumericalProtection} & Switch & \code{false} \\
\code{TNNPClampGates} & Switch & \code{true} \\
\code{TNNPClampVmForRates} & Switch & \code{true} \\
\code{TNNPLogNumericalProtection} & Switch & \code{true} \\
\code{TNNPConcentrationFloor} & scalar (mM) & $10^{-12}$ \\
\code{TNNPVMinForRates} & scalar (mV) & $-200$ \\
\code{TNNPVMaxForRates} & scalar (mV) & $+100$ \\
\code{TNNPVoltageZeroEps} & scalar (mV) & $10^{-8}$ \\
\code{TNNPMaxProtectionLogEntries} & label & $10^{6}$ \\
\bottomrule
\end{tabular}
\end{center}

When protection is enabled, the guards applied in order are:
\begin{enumerate}[topsep=0pt,itemsep=2pt]
  \item \textbf{Non-finite $V$}: \code{!std::isfinite(V)} $\rightarrow$
        reset to $-86.2$ mV.
  \item \textbf{$V$ clamp for rates} (if \code{TNNPClampVmForRates}):
        clamp to $[\,$\code{TNNPVMinForRates}, \code{TNNPVMaxForRates}$\,]$.
  \item \textbf{$V=0$ singularity avoidance}: push $V$ off zero by
        \code{TNNPVoltageZeroEps} (sign preserving) when
        $\lvert V \rvert$ falls below that threshold.
  \item \textbf{Concentration floor} on $K_{i}, Na_{i}, Ca_{i}, Ca_{SR}$:
        non-finite or below-floor entries are reset to
        \code{TNNPConcentrationFloor}.
  \item \textbf{Gate clamp} on $Xr_{1}, Xr_{2}, Xs, m, h, j, d, f, fCa,
        s, r, g$ (if \code{TNNPClampGates}): non-finite to zero,
        clamp to $[0, 1]$.
\end{enumerate}

Each correction is bucketed by \code{(phase, variable, reason)} in an
in-memory \code{ProtectionSummaryEntry} table and emitted to
\code{TNNP\_numericalProtection\_summary.dat} at exit.

% ================================================================== %
\section{Activation tracking and benchmarks}
% ================================================================== %

\paragraph{Activation time per cell.}
First-crossing of $V_{\text{thr}}$ in a time step is recorded by linear
interpolation:
\begin{equation}
  w \;=\; \mathrm{clip}\!\Bigl(
      \frac{V_{\text{thr}} - V^{n}}{V^{n+1} - V^{n}},
      \; 0, 1
    \Bigr),
   \qquad
   t_{\text{act}} \;=\; t^{n} + w\, \Delta t,
\end{equation}
with a per-cell flag cleared after the first crossing so subsequent
re-crossings are ignored.

\paragraph{Early termination on P8.}
When \code{stopAfterPointActivation = true}, the cell nearest
\code{stopActivationPoint} (P8 in the Niederer benchmark) is monitored
and \code{effectiveEndTime} is set to $t_{\text{act}} +
$~\code{stopDelayAfterActivation} once its $t_{\text{act}} > 0$.

\paragraph{Diagonal sampling.}
\code{sampleIDW} performs inverse-distance-weighted interpolation at
\code{nDiagonalSamples} points between \code{diagonalStart} and
\code{diagonalEnd}, with $k = $~\code{sampleNeighbours},
$p = $~\code{samplePower}.

% ================================================================== %
\section{Solver flow chart}
% ================================================================== %

The pseudo-code below summarises the per-time-step flow of the solver
and shows the \pyconst{run\_convergence.py constants}~/ \dictkey{dict
keys} that control each branch.

\begin{algorithm}[H]
\caption{High-level time-step loop of
         \texttt{highOrderElectroActivationFoamImplicitJfNK}.}
\begin{algorithmic}[1]
\State Initialise PETSc-free Eigen back-end and read dictionaries.
\State Assemble $\Mmass$ (\dictkey{massMatrix}), $\Kstif$
       (\dictkey{useHighOrder\_Vm}), $\Lop = \Kstif/(\chi C_{m})$.
\State $\mathsf{V}^{0} \gets -84$~mV (initial condition).
\State $\mathsf{V}^{-1} \gets \mathsf{V}^{0}$,\ \,
       $\Delta t_{\text{prev}} \gets 0$,\ \,
       \texttt{hasPrev} $\gets$ false.
\While{ $t^{n} < $~\pyconst{END\_TIME} }
   \State $\Delta t \gets $~\pyconst{DT\_VALUES} (current sweep value).
   \State Build $\Aimp = \Mmass/\Delta t - \theta\,\Lop$,
          $\Bimp = \Mmass/\Delta t + (1-\theta)\,\Lop$
          (\dictkey{implicitScheme} $\to \theta$).
   \State Apply stimulus to $\Iext$ (\dictkey{stimulusProtocol}).
   \If{ \dictkey{jfnkInitGuessOrder} $\ge 1$ \textbf{and} \texttt{hasPrev} }
      \State $\mathsf{V}_{\text{guess}} \gets \mathsf{V}^{n}
             + (\Delta t/\Delta t_{\text{prev}})\,(\mathsf{V}^{n} - \mathsf{V}^{n-1})$,
             clamp to $[\,$\dictkey{jfnkInitGuessVmMin},
             \dictkey{jfnkInitGuessVmMax}$\,]$.
   \Else
      \State $\mathsf{V}_{\text{guess}} \gets \mathsf{V}^{n}$.
   \EndIf
   \State $\mathsf{V}^{n-1} \gets \mathsf{V}^{n}$,\ \,
          $\Delta t_{\text{prev}} \gets \Delta t$,\ \,
          \texttt{hasPrev} $\gets$ true.
   \If{ \dictkey{nonlinearMethod} = JFNK }
      \State $\mathsf{x} \gets \mathsf{V}_{\text{guess}}$.
      \For{ $k = 0, 1, \dots, $~\pyconst{IMPLICIT\_NONLINEAR\_ITERATIONS}$-1$ }
         \State $\mathsf{F}_{k} \gets$ \Call{computeResidual}{$\mathsf{x}$}
                (advances TNNP via \dictkey{ODE\_REL\_TOL}, \dictkey{ODE\_ABS\_TOL},
                clamps Vm by \dictkey{jfnkClampODEInput}).
         \If{ all convergence tests pass (\dictkey{nonlinearTolerance}, $\dots$) }
            \State \textbf{break}
         \EndIf
         \State Build matrix-free Jacobian:
                $\mathsf{J}\mathbf{v} \approx (\mathsf{F}(\mathsf{x}+\varepsilon\mathbf{v})-\mathsf{F}_{k})/\varepsilon$
                with $\varepsilon$ from \dictkey{jfnkEpsilon}.
         \State Solve $\mathsf{J}\,\boldsymbol{\delta} = -\mathsf{F}_{k}$ with
                in-house GMRES (\dictkey{jfnkMaxKrylovIterations},
                \dictkey{jfnkLinearTolerance}).
         \If{ \dictkey{jfnkLineSearch} }
            \State Armijo backtracking (\dictkey{jfnkArmijoC},
                   \dictkey{jfnkLineSearchMaxIter},
                   \dictkey{jfnkLineSearchAlphaMin}):
                   find $\alpha$ satisfying
                   $\lVert\mathsf{F}(\mathsf{x}+\alpha\boldsymbol{\delta})\rVert^{2}
                    \le (1-2c\alpha)\lVert\mathsf{F}_{k}\rVert^{2}$.
         \Else
            \State $\alpha \gets $ \dictkey{nonlinearRelaxation}.
         \EndIf
         \State $\mathsf{x} \gets \mathsf{x} + \alpha\,\boldsymbol{\delta}$.
      \EndFor
      \State $\mathsf{V}^{n+1} \gets \mathsf{x}$.
   \ElsIf{ \dictkey{nonlinearMethod} = picard or diagonalIion }
      \For{ $k = 0, 1, \dots, $~\pyconst{IMPLICIT\_NONLINEAR\_ITERATIONS}$-1$ }
         \State Solve $\Aimp\,\mathsf{V}_{*} = \mathsf{r}(\Iion(\mathsf{x}_{k}))$
                with Eigen (\dictkey{linearSolver},
                \dictkey{tolerance}).
         \State $\mathsf{x}_{k+1} \gets
                (1{-}\omega)\,\mathsf{x}_{k} + \omega\,\mathsf{V}_{*}$,
                $\omega = $~\dictkey{nonlinearRelaxation}.
         \If{ convergence tests pass } \textbf{break} \EndIf
      \EndFor
      \State $\mathsf{V}^{n+1} \gets \mathsf{x}_{k+1}$.
   \EndIf
   \If{ \textbf{not} converged \textbf{and not} \dictkey{nonlinearAcceptUnconverged} }
      \State Warn; roll back: $\mathsf{V}^{n+1} \gets \mathsf{V}^{n}$,
             $\Iion^{n+1} \gets \Iion^{n}$, $\mathbf{s}^{n+1} \gets \mathbf{s}^{n}$.
   \EndIf
   \State Record activation times (\dictkey{activationThreshold}),
          P8 stop check (\dictkey{stopAfterPointActivation},
          \dictkey{stopDelayAfterActivation}).
   \State Write nonlinear-residual log if \dictkey{writeNonlinearResiduals}.
   \State $t^{n+1} \gets t^{n} + \Delta t$.
\EndWhile
\State Post-process: IDW sampling at P8 and diagonal
       (\dictkey{nDiagonalSamples}, \dictkey{sampleNeighbours},
        \dictkey{samplePower}).
\end{algorithmic}
\end{algorithm}

\noindent\textbf{Legend.} \pyconst{Blue} = Python constant set in
\code{run\_convergence.py}; \dictkey{green} = OpenFOAM dictionary key
read by the solver. Constants and keys are mapped one-to-one in
section~\ref{sec:driver}.

% ================================================================== %
\section{Driver script: \texttt{run\_convergence.py}}
\label{sec:driver}
% ================================================================== %

The Python driver generates the OpenFOAM dictionaries and runs the
solver in a parameter sweep. Tables below group constants by topic;
\emph{Key} is the dictionary entry emitted by the script.

\subsection{Geometry and parameter sweep}
\begin{longtable}{@{}p{0.30\linewidth}p{0.28\linewidth}p{0.36\linewidth}@{}}
\toprule
Python constant & OpenFOAM key & Meaning \\
\midrule \endhead
\code{DIMENSIONS} & --- &
   List of dimensions to sweep (\code{2D}, \code{3D}). \\
\code{DX\_VALUES\_MM} & \code{blockMeshDict} cell counts &
   Uniform cell size in millimetres. \\
\code{DT\_VALUES} & \code{controlDict.deltaT} &
   Outer time-step sizes (seconds). \\
\code{END\_TIME} & \code{controlDict.endTime} & End time (s). \\
\code{CFL} & \code{timeIntegrationProperties.CFL} &
   Informational reference; not enforced. \\
\code{ACTIVE\_TIME\_SCHEMES} &
   \code{spatialIntegrationProperties.implicitScheme} &
   \{\code{backwardEuler}, \code{crankNicolson}\}. \\
\code{ACTIVE\_MASS\_MATRICES} &
   \code{spatialIntegrationProperties.massMatrix} &
   \{\code{lumped}, \code{consistent}\}. \\
\code{ACTIVE\_HO\_PAIRS\_Vm\_States} &
   \code{useHighOrder\_Vm} / \code{useHighOrder\_Iion}
   + \code{LRECoeffs\_*} &
   Tuples \texttt{(Vm-level, Iion-level)}; level $\in$
   \{\code{NO}, \code{p1}, \code{p2}, \code{p3}\}. \\
\code{ACTIVE\_NONLINEAR\_METHODS} &
   \code{nonlinearMethod} &
   \{\code{JFNK}, \code{picard}, \code{diagonalIion}, $\dots$\}. \\
\code{ACTIVE\_STATE\_ODE\_SOLVERS} &
   \code{timeIntegrationProperties.solver} &
   ODE integrator names (\code{RKF45}, $\dots$). \\
\code{SIGMA\_X, SIGMA\_Y, SIGMA\_Z} & \code{cardiacProperties} &
   Diagonal of $\Dten$. \\
\code{CHI} & \code{cardiacProperties.chi} & Surface-to-volume ratio. \\
\code{CM} & \code{cardiacProperties.Cm} & Membrane capacitance. \\
\bottomrule
\end{longtable}

\subsection{Nonlinear solver (JFNK / Picard)}
\begin{longtable}{@{}p{0.35\linewidth}p{0.30\linewidth}p{0.30\linewidth}@{}}
\toprule
Python constant & OpenFOAM key & Meaning \\
\midrule \endhead
\code{IMPLICIT\_NONLINEAR\_ITERATIONS} & \code{nonlinearIterations} &
   Newton/Picard iteration cap per step. \\
\code{IMPLICIT\_MIN\_NONLINEAR\_ITERATIONS} &
   \code{minNonlinearIterations} & Minimum Newton iterations. \\
\code{NONLINEAR\_TOLERANCE} & \code{nonlinearTolerance} &
   Coupled-residual tolerance
   $\lVert\mathsf{F}\rVert/\lVert\mathsf{x}\rVert$. \\
\code{NONLINEAR\_VM\_TOLERANCE} & \code{nonlinearVmTolerance} &
   Relative L$^{2}$ change in $\Vm$. \\
\code{NONLINEAR\_IION\_TOLERANCE} & \code{nonlinearIionTolerance} &
   Same on $\Iion$. \\
\code{NONLINEAR\_STATES\_TOLERANCE} &
   \code{nonlinearStatesTolerance} & Max over states. \\
\code{NONLINEAR\_REQUIRE\_STATES\_CONVERGENCE} &
   \code{nonlinearRequireStatesConvergence} & Gate state test. \\
\code{NONLINEAR\_RELAXATION} & \code{nonlinearRelaxation} &
   $\alpha_{0}$ (Newton damping / Picard mix). \\
\bottomrule
\end{longtable}

\subsection{JFNK matrix-free options}
\begin{longtable}{@{}p{0.35\linewidth}p{0.30\linewidth}p{0.30\linewidth}@{}}
\toprule
Python constant & OpenFOAM key & Meaning \\
\midrule \endhead
\code{JFNK\_MAX\_KRYLOV\_ITERATIONS} & \code{jfnkMaxKrylovIterations} &
   GMRES restart cycle $m$. \\
\code{JFNK\_LINEAR\_TOLERANCE} & \code{jfnkLinearTolerance} &
   GMRES relative tolerance. \\
\code{JFNK\_EPSILON} & \code{jfnkEpsilon} &
   FD perturbation $\varepsilon_{0}$. Set near $\sqrt{\eta}$ with
   $\eta = $~\code{relTol}. \\
\code{DIAGONAL\_IION\_EPSILON} & \code{diagonalIionEpsilon} &
   FD step for diagonal $\partial \Iion/\partial \Vm$. \\
\bottomrule
\end{longtable}

\subsection{Linear solver (Eigen)}
\begin{longtable}{@{}p{0.30\linewidth}p{0.30\linewidth}p{0.35\linewidth}@{}}
\toprule
Python constant & OpenFOAM key & Meaning \\
\midrule \endhead
\code{IMPLICIT\_LINEAR\_SOLVER} & \code{linearSolver} &
   \{\code{SparseLU}, \code{BiCGSTAB}\}. Used by Picard / diagonal
   branches only; JFNK uses the in-house GMRES on the FD shell. \\
\code{IMPLICIT\_TOLERANCE} & \code{tolerance} &
   BiCGSTAB relative tolerance. \\
\code{IMPLICIT\_MAX\_ITERATIONS} & \code{maxIterations} &
   BiCGSTAB iteration cap. \\
\bottomrule
\end{longtable}

\subsection{Inner ODE integrator (RKF45)}
\begin{longtable}{@{}p{0.30\linewidth}p{0.30\linewidth}p{0.35\linewidth}@{}}
\toprule
Python constant & OpenFOAM key & Meaning \\
\midrule \endhead
\code{ODE\_ABS\_TOL} & \code{timeIntegrationProperties.absTol} &
   Absolute tolerance per state. \\
\code{ODE\_REL\_TOL} & \code{timeIntegrationProperties.relTol} &
   Relative tolerance per state; drives the F-noise floor. \\
\code{ODE\_INITIAL\_STEP} & \code{initialODEStep} &
   Initial sub-step (ms). \\
\code{ODE\_MAX\_STEPS} & \code{maxSteps} &
   Cap on RKF45 sub-steps per outer call. \\
\bottomrule
\end{longtable}

\subsection{LRE high-order parameters per p-level}
\code{P\_LEVELS\_BY\_DIM[dim][p]} gives, for each (dim, p):
\code{N} (polynomial order), \code{Nn} (stencil neighbours),
\code{k} (search radius), \code{maxStencilSize} (cap). The script
writes these into a nested
\code{highOrderCoeffs\{LRECoeffs\_Vm\{\dots\}; LRECoeffs\_Iion\{\dots\}\}}
block whenever HO is enabled for either Vm or $\Iion$.

\subsection{Post-processing and run-control}
\begin{longtable}{@{}p{0.35\linewidth}p{0.30\linewidth}p{0.30\linewidth}@{}}
\toprule
Python constant & OpenFOAM key & Meaning \\
\midrule \endhead
\code{ACTIVATION\_THRESHOLD} & \code{activationThreshold} &
   $V_{\text{thr}}$ (V) for activation-time crossing. \\
\code{N\_DIAGONAL\_SAMPLES} & \code{nDiagonalSamples} &
   IDW samples along the diagonal. \\
\code{SAMPLE\_NEIGHBOURS} & \code{sampleNeighbours} & $k$-NN for IDW. \\
\code{SAMPLE\_POWER} & \code{samplePower} & Inverse-distance power. \\
\code{STOP\_AFTER\_POINT\_ACTIVATION} &
   \code{stopAfterPointActivation} &
   Early-stop on P8. \\
\code{STOP\_DELAY\_AFTER\_ACTIVATION} &
   \code{stopDelayAfterActivation} & Delay after P8 activation (s). \\
\code{WRITE\_NONLINEAR\_RESIDUALS} &
   \code{writeNonlinearResiduals} & Residual log on/off. \\
\code{SAVE\_LOGS\_ON\_SUCCESS} & --- &
   Python-only; keeps logs after a successful run. \\
\bottomrule
\end{longtable}

% ================================================================== %
\appendix
\section{Dimension table}
% ================================================================== %

\begin{center}
\begin{tabular}{@{}lll@{}}
\toprule
Field & OpenFOAM \code{dimensionSet} & SI \\
\midrule
\code{Vm} & $[1, 2, -3, 0, 0, -1, 0]$ & volt \\
\code{Iion} & \code{dimVoltage/dimTime} & \si{\volt\per\second} \\
\code{externalStimulusCurrent} & \code{dimCurrent/dimVolume} &
    \si{\ampere\per\meter\cubed} \\
\code{conductivity} &
    \code{dimCurrent\^{}2 dimTime\^{}3 / (dimMass dimVolume)} &
    \si{\siemens\per\meter} \\
\code{chi} & \code{dimless/dimLength} & \si{\per\meter} \\
\code{Cm} & \code{dimCurrent dimTime / (dimVoltage dimArea)} &
    \si{\farad\per\meter\squared} \\
\bottomrule
\end{tabular}
\end{center}

\section{File map}

\begin{center}
\begin{tabular}{@{}p{0.55\linewidth}p{0.40\linewidth}@{}}
\toprule
Path (relative to \code{cardiacFoam-pc/}) & Role \\
\midrule
\code{applications/solvers/.../JfNK.C} & Main translation unit. \\
\code{applications/solvers/.../createFields.H} & Field initialisation
   and dictionary reading. \\
\code{src/ionicModels/TNNP/TNNP.\{H,C\}} & TNNP cell model and
   protection guards. \\
\code{src/ionicModels/TNNP/TNNP\_2004.H, *Names.H} & Generated CellML
   kernel. \\
\code{src/ionicModels/ionicModel/ionicModel.\{H,C\}} & Base class
   with run-time ODE solver selection. \\
\code{tutorials\_electro/.../run\_convergence.py} & Driver script. \\
\code{constant/cardiacProperties} & $\chi$, $C_{m}$, $\Dten$,
   ionic model. \\
\code{constant/timeIntegrationProperties} & ODE solver, tolerances,
   sub-step cap. \\
\code{constant/spatialIntegrationProperties} & Time scheme, HO flags,
   JFNK / linear / line-search options. \\
\code{constant/stimulusProtocol} & Stimulus regions and timing. \\
\bottomrule
\end{tabular}
\end{center}

% ================================================================== %
\section{Differences vs.\ the PETSc sibling}
% ================================================================== %

Algorithmically the two solvers are identical. The only differences
are in the linear-algebra back-end and a few user-visible knobs:

\begin{center}
\begin{tabular}{@{}p{0.42\linewidth}p{0.26\linewidth}p{0.26\linewidth}@{}}
\toprule
Aspect & This solver (Eigen) & PETSc sibling \\
\midrule
Sparse matrix container & \code{Eigen::SparseMatrix} & PETSc \code{Mat} \\
Sparse direct solver & \code{Eigen::SparseLU} & \code{KSPPREONLY+PCLU} \\
Iterative solver (Picard branch) & \code{Eigen::BiCGSTAB+ILUT} &
    \code{KSP+PC} (full PETSc catalogue) \\
JFNK inner Krylov & in-house GMRES & \code{KSPGMRES} on \code{MatShell} \\
Preconditioner inside JFNK & none (Krylov-only) & any PETSc PC \\
Build-time dependency & header-only Eigen3 &
    PETSc + MPI \\
Parallel run support & serial only & MPI-parallel out of the box \\
\bottomrule
\end{tabular}
\end{center}

All convergence-related parameters (\code{nonlinearTolerance},
\code{jfnkEpsilon}, \code{jfnkLineSearch*}, \code{jfnkInitGuess*},
\code{jfnkClampODEInput}, \code{nonlinearAcceptUnconverged},
\code{ODE\_*}) have the same names and meanings in both solvers, so
the same \code{run\_convergence.py} configuration produces an
algorithmically equivalent run with either back-end (modulo the
performance and parallel-scaling differences).

% ================================================================== %
\begin{thebibliography}{9}
% ================================================================== %

\bibitem{ttp2004} K.~H.~W.~J. ten Tusscher, D. Noble, P.~J. Noble, A.~V.
    Panfilov, \emph{A model for human ventricular tissue},
    AJP -- Heart and Circulatory Physiology
    \textbf{286} (2004) H1573--H1589.

\bibitem{niederer2012} S.~A. Niederer et al.,
    \emph{Verification of cardiac tissue electrophysiology simulators
    using an N-version benchmark}, Prog.\ Biophys.\ Mol.\ Biol.
    \textbf{104} (2012) 22--48.

\bibitem{knoll2004} D.~A. Knoll, D.~E. Keyes,
    \emph{Jacobian-Free Newton-Krylov methods: a survey of approaches
    and applications}, J.\ Comp.\ Phys.\ \textbf{193}
    (2004) 357--397.

\bibitem{noceWright} J. Nocedal, S.~J. Wright,
    \emph{Numerical Optimization}, 2nd edition, Springer, 2006.

\bibitem{saadbook} Y. Saad,
    \emph{Iterative Methods for Sparse Linear Systems}, 2nd edition,
    SIAM, 2003.

\bibitem{hairer2} E. Hairer, G. Wanner,
    \emph{Solving Ordinary Differential Equations II}, Springer, 1996.

\end{thebibliography}

\end{document}
